\section{Introducción}
El desarrollo de software ha experimentado transformaciones significativas en las últimas décadas. La creciente complejidad de los sistemas exige soluciones más organizadas, escalables y seguras. En este escenario, los paradigmas de programación, las metodologías de gestión y las arquitecturas de software se convierten en piezas fundamentales para enfrentar los retos tecnológicos actuales. Según \cite{scrum}, la ausencia de una arquitectura sólida impacta negativamente en el ciclo de vida del software, dificultando su evolución y mantenibilidad. De igual manera, Gavilánez et al. \cite{poo} destacan que los patrones de diseño y la POO garantizan la claridad del código y su reutilización. Por su parte, \cite{nodejs_springboot} afirman que la adopción de metodologías ágiles como Scrum permite mejorar la comunicación, adaptarse al cambio y entregar valor continuo. En ese sentido, la integración de metodologías y arquitecturas ofrece un marco sólido para la creación de sistemas más eficientes. Este artículo busca integrar los aportes teóricos y prácticos de la POO, el patrón MVC, los frameworks Node.js y Spring Boot, la metodología Scrum y la gestión de bases de datos, con el fin de mostrar su relevancia en el desarrollo de software moderno.
